%% GENERATED by experiments/019-simb-multimodal/scripts/morph_v24_dataset_table.py. Do not edit by hand.
%% SOURCE: the two stores' processed/phenotype_label_index.json and data_module_cache/index_seed_0.json;
%% results/morphology_noise_ceiling.json, morphology_feature_ceiling.csv, calmorph_moving_features.json,
%% proteome_morphology_covariation.json, expression_morphology_covariation.json, morphology_baselines_split_fig3_core/
\begin{table}[htbp]
  \centering
  \footnotesize
  \caption[The morphology round's strain sets]{The morphology round's strain sets. Genotype records per store and per label, the overlaps, the split seed 0 partition, and per arm the rows left after the training script's \texttt{require\_modalities} filter (the counts the runs print). A genotype is one deletion strain after deduplication and aggregation; a strain measured by two panels is one record carrying both labels. \src{experiments/019-simb-multimodal/scripts/morph_v24_dataset_table.py}}
  \label{tab:morph-strains}
  \begin{tabular}{@{}llrrrrrrr@{}}
    \toprule
    store & sources & genotypes & morphology & expression & proteome & morph.\ and expr.\ & morph.\ and prot.\ & all three \\
    \midrule
    \texttt{fig3\_core} & Kemmeren 2014, Sameith 2015, Ohya 2005, Costanzo 2016 and Kuzmin 2018/2020 fitness & 5,125 & 4,695 & 1,554 & -- & 1,439 & -- & -- \\
    \texttt{fig3\_morph\_proteome} & Kemmeren 2014, Sameith 2015, Messner 2023, Ohya 2005 & 4,977 & 4,695 & 1,554 & 4476 & 1,439 & 4281 & 1321 \\
    \midrule
    \multicolumn{9}{@{}l}{\emph{Per arm, split seed 0, after the label filter (train / validation / test rows)}}\\
    \multicolumn{2}{@{}l}{\texttt{M\_cls, M\_pool} on \texttt{fig3\_core}} & \multicolumn{7}{l}{morphology: 3,757 / 469 / 469} \\
    \multicolumn{2}{@{}l}{\texttt{CM\_expr, CM\_exprperm, CMP\_expr, CMP\_exprperm} on \texttt{fig3\_core}} & \multicolumn{7}{l}{expression, morphology: 1,161 / 132 / 146} \\
    \multicolumn{2}{@{}l}{\texttt{MP\_cls, CM\_prot, CM\_protperm} on \texttt{fig3\_morph\_proteome}} & \multicolumn{7}{l}{proteome, morphology: 3,464 / 416 / 401} \\
    \bottomrule
  \end{tabular}
\end{table}

\begin{table}[htbp]
  \centering
  \footnotesize
  \caption[The morphology target]{The morphology target and its references. CalMorph (Ohya 2005) gives one population-averaged profile per deletion and 122 wild-type replicates; a feature's reliability is one minus the replicate variance over the variance across mutants, its ceiling the square root, the best across-strain Pearson a perfect predictor of the signal can reach against one noisy measurement. The round predicts the moving features only. The ridge rows are the out-of-fold median Pearson per feature from the measured other modality on the strains both panels measured; the sequence rows are the best key selected on validation, four split seeds, all 278 features. \src{experiments/019-simb-multimodal/scripts/morph_v24_dataset_table.py}}
  \label{tab:morph-target}
  \begin{tabular}{@{}lr@{}}
    \toprule
    quantity & value \\
    \midrule
    deletion strains measured (Ohya 2005) & 4,718 \\
    wild-type replicates & 122 \\
    base CalMorph features & 281 \\
    after dropping the three degenerate & 278 \\
    moving features, reliability $\geq$ 0.5 (the target) & 116 \\
    target transform & Yeo-Johnson per feature, fit on the training rows \\
    ceiling, mean over all 278 features & 0.611 \\
    ceiling, mean over the moving features (minimum) & 0.847 (0.712) \\
    ridge from the measured proteome, moving features & 0.280 \\
    ridge from measured expression, moving features & 0.324 \\
    sequence bilinear ridge, all 278 features & 0.086 \\
    sequence neighbor average, all 278 features & 0.088 \\
    \bottomrule
  \end{tabular}
\end{table}
